\documentclass[11pt,letterpaper]{article}
\usepackage[lmargin=1in,rmargin=1in,tmargin=1in,bmargin=1in]{geometry}

% -------------------
% Packages
% -------------------
\usepackage{
	amsmath,			% Math Environments
	amssymb,			% Extended Symbols
	amsthm,			    % Theorem Environments
	cancel,			    % Use Cancels
	enumerate,		    % Enumerate Environments
	graphicx,			% Include Images
	lastpage,			% Reference Lastpage
	multicol,			% Use Multi-columns
	multirow,			% Use Multi-rows
	xcolor			    % Use Colors
}


% -------------------
% Font
% -------------------
\usepackage[T1]{fontenc}
\usepackage{charter}

%\usepackage[T1]{fontenc}
%\usepackage{mathpazo}

%\usepackage[bitstream-charter]{mathdesign}
%\usepackage[T1]{fontenc}


% -------------------
% Tikz & PGF
% -------------------
\usepackage{tikz}
\usepackage{tikz-cd}
\usetikzlibrary{
	calc,
	decorations.pathmorphing,
	matrix,arrows,
	positioning,
	shapes.geometric
}
\usepackage{pgfplots}
\pgfplotsset{compat=newest}


% -------------------
% Commands
% -------------------

% Problem Labels
\newcounter{problem}
\newcommand{\problem}{
\stepcounter{problem} %
\noindent \textbf{Problem \theproblem. }}
\newcommand{\solution}{\noindent \textbf{Solution: }}
\newcommand{\pf}{\noindent\emph{Proof. }}

% Special Characters
\newcommand{\N}{\mathbb{N}}
\newcommand{\Z}{\mathbb{Z}}
\newcommand{\Q}{\mathbb{Q}}
\newcommand{\R}{\mathbb{R}}
\newcommand{\C}{\mathbb{C}}

% Math Operators
\DeclareMathOperator{\im}{im}
\DeclareMathOperator{\disc}{disc}
\DeclareMathOperator{\Span}{span}

% Special Commands
\newcommand{\pspace}{\par\vspace{\baselineskip}}
\newcommand{\ds}{\displaystyle}
\newcommand{\defeq}{\stackrel{\text{def}}{=}}
\newcommand{\ov}[1]{\overline{#1}}
\newcommand{\ma}[1]{\stackrel{#1}{\longrightarrow}}
\newcommand{\twomatrix}[4]{\begin{pmatrix} #1 & #2 \\ #3 & #4 \end{pmatrix}}


% -------------------
% Header & Footer
% -------------------
\usepackage{fancyhdr}

\fancypagestyle{title}{
	%Headers
	\fancyhead[L]{\Large My Name}
	\fancyhead[C]{\Large MAT 999}
	\fancyhead[R]{\Large Homework \#}
	\renewcommand{\headrulewidth}{0.2pt}
	%Footers
	\fancyfoot[L]{}
	\fancyfoot[C]{}
	\fancyfoot[R]{\thepage \,of \pageref{LastPage}}
	\renewcommand{\footrulewidth}{0.0pt}
}
\fancypagestyle{pages}{
	%Headers
	\fancyhead[L]{}
	\fancyhead[C]{}
	\fancyhead[R]{}
	\renewcommand{\headrulewidth}{0.0pt}
	%Footers
	\fancyfoot[L]{}
	\fancyfoot[C]{}
	\fancyfoot[R]{\thepage \,of \pageref{LastPage}}
	\renewcommand{\footrulewidth}{0.0pt}
}
\headheight=18pt
\footskip=14pt

\pagestyle{pages}


% -------------------
% Content
% -------------------
\begin{document}
\thispagestyle{title}



% Problem 1
\problem Show that there exists no nontrivial unramified extensions of $\Q$. \pspace

\solution  If $K/\Q$ is a nontrivial number field, then $|\disc K|>1$. But then $\disc K$ has a prime factor so that some prime ramifies in $K$. \qed \pspace



% Problem 2
\problem Complete the following:
        \begin{enumerate}[(a)]
	\item How does one prove a cotheorem?
	\item Compute $\ds \int \cos x \;dx$.
	\item How does one square $\twomatrix{a}{b}{c}{d}$?       
        \end{enumerate}

\solution
\begin{enumerate}[(a)]
\item Use rollaries.

\item We have
	\begin{equation} \label{eq:integral}
	\int \cos x \;dx= \sin x + C	
	\end{equation}
We can check \eqref{eq:integral}:
	\[
	\dfrac{d}{dx} \left( \sin x + C \right)= \cos x
	\]

\item This is routine. 
\end{enumerate} \qed \pspace



% Problem 3
\problem Prove that $\sqrt{2}$ is irrational. \pspace

\pf Assume that $\sqrt{2}= \dfrac{a}{b}$, where $a,b \in \Z$. Without loss of generality, we may assume $\gcd(a,b)= 1$. Then we have
	\begin{align}
	\sqrt{2}&= \dfrac{a}{b} \nonumber \\
	\sqrt{2}^2&= \left( \dfrac{a}{b} \right)^2 \label{eq:implication1} \\
	2&= \dfrac{a^2}{b^2} \nonumber \\
	a^2&= 2b^2 \label{eq:implication2}
	\end{align}
But then from \eqref{eq:implication2}, we know that $a^2$ is even so that $a$ is even. But then we must have
	\[
	2a^2= b^2
	\]
so that $b^2$ is even, implying $b$ is even. But then $\gcd(a,b) \geq 2$, a contradiction. \qed \\



\end{document}